% The three assembly chemistries FORGE generates into, drawn natively so every glyph is the
% document's own font.  Colours follow the role palette of the algorithm schematic: the amine head
% is roleamine, aldehyde-derived atoms roleald, the isocyanide or acrylate tail roleiso, and an atom
% the template introduces with no precursor origin is grey, exactly as in the Figure 1 core inset.
% Skeletal notation throughout; hydrogens are drawn only on the atoms whose substitution decides
% whether the product can be decomposed back into its precursors.  Bonds carry the colour of the
% component that contributes them, so the reader can read the product back to its three sources.
\definecolor{roleamine}{HTML}{8A5AC2}
\definecolor{roleald}{HTML}{168C8C}
\definecolor{roleiso}{HTML}{E39532}
\definecolor{roletemplate}{HTML}{9A9A9A}
\setchemfig{atom sep=2.35em, bond style={line width=0.9pt}}
\newcommand{\amine}[1]{{\color{roleamine}#1}}
\newcommand{\ald}[1]{{\color{roleald}#1}}
\newcommand{\iso}[1]{{\color{roleiso}#1}}
\newcommand{\tmpl}[1]{{\color{roletemplate}#1}}
\newcommand{\rxnrow}[3]{%
  \makebox[\linewidth][l]{\begin{minipage}{\linewidth}%
  \textbf{\small #1}\\[-0.15em]{\scriptsize\itshape #2}\\[0.7em]#3\end{minipage}}}

\rxnrow{Ugi three-component reaction}
       {amine head, aldehyde and isocyanide combine in one step; no carboxylic-acid reactant component}{%
\schemestart
  \chemfig{R^1-[:30]\amine{N}(-[:90]R^2)-[:-30]\amine{H}}
  \+
  \chemfig{R^3-[:30,,,,roleald]\ald{CH}=[:-30,,,,roleald]\ald{O}}
  \+
  \chemfig{R^4-[:30,,,,roleiso]\iso{N}~[:-30,,,,roleiso]\iso{C}}
  \arrow{->}
  \chemfig{R^1-[:30,,,,roleamine]\amine{N}(-[:90,,,,roleamine]R^2)-[:-30,,,,roleald]\ald{CH}(-[:-90,,,,roleald]R^3)-[:30,,,,roleiso]\iso{C}(=[:90,,,,roletemplate]\tmpl{O})-[:-30,,,,roleiso]\iso{NH}-[:30,,,,roleiso]R^4}
\schemestop}\par

\par\vspace{1.9em}\par
\rxnrow{Repeated aza-Michael addition}
       {the head N--H adds across the acrylate double bond; the ester is untouched}{%
\schemestart
  \chemfig{R^1-[:30]\amine{N}(-[:90]R^2)-[:-30]\amine{H}}
  \+
  \chemfig{-[:30,,,,roleiso]=[:-30,,,,roleiso]-[:30](=[:90]O)-[:-30]O-[:30]R^3}
  \arrow{->}
  \chemfig{R^1-[:30,,,,roleamine]\amine{N}(-[:90,,,,roleamine]R^2)-[:-30,,,,roleiso]-[:30,,,,roleiso]-[:-30](=[:90]O)-[:30]O-[:-30]R^3}
\schemestop}\par

\par\vspace{1.9em}\par
\rxnrow{Repeated reductive amination}
       {one new carbon--nitrogen bond; the carbonyl oxygen leaves as water}{%
\schemestart
  \chemfig{R^1-[:30]\amine{N}(-[:90]R^2)-[:-30]\amine{H}}
  \+
  \chemfig{R^3-[:30,,,,roleald]\ald{CH}=[:-30,,,,roleald]\ald{O}}
  \arrow{->[\scriptsize$-\,$H$_2$O]}
  \chemfig{R^1-[:30,,,,roleamine]\amine{N}(-[:90,,,,roleamine]R^2)-[:-30,,,,roleald]-[:30,,,,roleald]R^3}
\schemestop}
